%%%%%%%%%%%%%%%%%%%%%%%%%
% Imports and Setup     %
%%%%%%%%%%%%%%%%%%%%%%%%%

\input{../lib/hw-setup/tex/HWSetup}
\input{../lib/hw-setup/tex/EngBindings}

\fancyhead[L]{\textbf{I.\ Dutta,\ M.\ Kostrzewa,\ \&\ V.\ Srivastava}}
\fancyhead[C]{}
\fancyhead[R]{\hyperlink{tableofcontents}{\hmwkClass\ (\hmwkClassInstructor,\ \hmwkClassTime): \hmwkTitle}}
\fancyheadwidth[L]{0.5\headwidth}
\fancyheadwidth[C]{0.0\headwidth}
\fancyheadwidth[R]{0.5\headwidth}

%%%%%%%%%%%%%%%%%%%%%%%%%
% Homework Details      %
% - Title               %
% - Subtitle            %
% - Due Date            %
% - Due Time            %
% - Course              %
% - Section Time        %
% - Instructor          %
% - Author              %
% - Submission Time     %
%%%%%%%%%%%%%%%%%%%%%%%%%

\hwkTitle{Lab03}
\hwkSubTitle{Choking in Convergent-Divergent Duct}
\hwkDueDate{2026-03-13}
\hwkDueTime{23:59:00}
\hwkClass{ENAE464 - 0202}
\hwkClassTime{14:00:00}
\hwkInstructor{Dr. Silbaugh}
\hwkAuthor{Ishan Dutta, Mikołaj Kostrzewa, \& Vai Srivastava}
\hwkCompletionDate{
        Experiment Performed: \DTMdate{2026-03-06}\\
        Report Submitted: \DTMtoday
}

\addbibresource{./references.bib}

\begin{document}

\hypertarget{titlepage}{\maketitle}
\thispagestyle{empty}

\pagebreak

%%%%%%%%%%%%%%%%%%%%%%%%%
% Letter of Transmittal %
%%%%%%%%%%%%%%%%%%%%%%%%%

\thispagestyle{empty}

\DTMtoday

\vspace{5.0ex}

Enclosed is the technical report for the \textit{Choking in Convergent-Divergent Duct} laboratory experiment. This report presents the experimental methodology, results, and analysis of mass flow rate for subsonic \& supersonic flow within a convergent-divergent duct.

Please feel free to contact us with any questions regarding the contents of this report.

\vspace{5.0ex}

Respectfully,

Ishan Dutta, Mikołaj Kostrzewa, \& Vai Srivastava

%%%%%%%%%%%%%%%%%%%%%%%%%
% Table of Contents     %
%%%%%%%%%%%%%%%%%%%%%%%%%

\pagebreak

\thispagestyle{empty}
\hypertarget{tableofcontents}{\tableofcontents}

\renewcommand{\listfigurename}{Figures}
\listoffigures

\renewcommand{\listtablename}{Tables}
\listoftables

\renewcommand{\listlistingname}{Listings}
\listoflistings

\newpage
\pagenumbering{arabic}

%%%%%%%%%%%%%%%%%%%%%%%%%
% Report                %
%%%%%%%%%%%%%%%%%%%%%%%%%

% FIXME: wip
\section{Abstract} \label{sec:abstract}

% TODO: abstract
\lipsum[1-2]

% FIXME: wip
\section{Introduction} \label{sec:introduction}

% TODO: basic theory and objectives
\lipsum[1-2]

% FIXME: wip
\section{Methodology} \label{sec:methodology}

% FIXME: wip
\subsection{Experimental Apparatus} \label{sec:methodology/apparatus}

% TODO: description and image of test setup
\lipsum[1-2]

% FIXME: wip
\subsection{Operating Procedure} \label{sec:methodology/procedure}

% TODO: methodology/procedure
\lipsum[1-2]

% FIXME: wip
\section{Results} \label{sec:results}

\subsection{Operating Conditions} \label{sec:results/conditions}

The density of air (\( \rho \)) can be calculated using the Ideal Gas Law:
\begin{equation}
    \rho = \frac{P}{R T}
\end{equation}
where:
\begin{itemize}
    \item \( P \) is the absolute pressure (in SI units: Pascals \unit{\Pa})
    \item \( T \) is the absolute temperature (in Kelvin \unit{\K})
    \item \( R \) is the specific gas constant for dry air, \( R = \qty{287.05}{\J\per\kg\per\K} \)
\end{itemize}

Given Laboratory Condition Measurements from Table \ref{tab:data/ambient-conditions}:
\begin{align}
    P &= \qty{1025.1}{\hecto\Pa} = \qty{102510}{\Pa} \\
    T &= \qty{22}{\C} = 22 + 273.15 = \qty{295.15}{\K} \\
    R &= \qty{287.05}{\J\per\kg\per\K}
\end{align}

\begin{equation}
    \rho = \frac{\qty{102510}{\Pa}}{\qty{287.05}{\J\per\kg\per\K}} \times \qty{295.15}{\K}
\end{equation}

Calculate:
\begin{align}
    \rho &= \frac{102510}{287.05 \times 295.15} \\
         &\approx \qty{1.20995}{\kg\per\m\cubed}
\end{align}

Thus, \textbf{the ambient air density during the experiment is}:
\begin{equation}
    \boxed{\rho = \qty{1.21}{\kg\per\m\cubed}}
\end{equation}

% FIXME: wip
\subsection{Measurements} \label{sec:results/measurements}

% TODO: table of relative pressures
% - atmosphere (reference)
% - inlet
% - throat
% - outlet
%
% % sample table import
% \begin{table}[H]
%     \centering
%     \caption{Measured Relative Presssure, Pressure Coefficient, and Pressure vs. Angle of Attack}
%     \mintinline{bash}{./outputs/text/pressure_vs_theta.csv} \\
%     \begin{tabular}[c]{|r|l|c|c|c|}\hline
%         & \( \theta \, \left[\unit{\deg}\right] \) & \( P-P_{\infty} \, \left[\unit{\Pa}\right] \) & \( C_{P} \) & \( P \, \left[\unit{\Pa}\right] \) \\ \hline \hline
%         \csvreader[late after line = \\ \hline]{../outputs/text/pressure_vs_theta.csv}{}{
%             {\setpadnum{2}\tt\scriptsize\padnum{\thecsvrow}} &
%             \csvcoli &
%             \csvcolii &
%             \csvcoliii &
%             \csvcoliv
%         }
%     \end{tabular}
%     \label{tab:results/measurements/pressures}
% \end{table}

% FIXME: wip
\subsection{Mass Flow Rate} \label{sec:results/mass-flow-rate}

% TODO: calculation of mass flow rate from pressures
\lipsum[1-2]

% FIXME: wip
\section{Analysis and Discussion} \label{sec:analysis}

% FIXME: wip
\subsection{Relations between Mass Flow Rate and Pressures} \label{sec:analysis/relations}

% TODO: plot of \( \dot{m} \) vs. \( p_{atmos} - p_{outlet} \)
%
% % sample figure import
% \begin{figure}[H]
%     \centering
%     \includegraphics[width=0.65\linewidth]{../outputs/figures/Cp_vs_theta.png}
%     \caption{Pressure Coefficient \( C_{p} \) vs. Angle of Attack \( \theta \)}
%     \label{fig:analysis/relations/Cp}
% \end{figure}

% TODO: plot of \( p_{atmos} - p_{throat} \) vs. \( p_{atmos} - p_{outlet} \)
%
% % sample figure import
% \begin{figure}[H]
%     \centering
%     \includegraphics[width=0.65\linewidth]{../outputs/figures/Cd_vs_theta.png}
%     \caption{Drag Coefficient \( C_{d} \) vs. Angle of Attack \( \theta \)}
%     \label{fig:analysis/relations/Cd}
% \end{figure}

% FIXME: wip
\subsection{Interpretation} \label{sec:analysis/interpretation}

% TODO: comment on the shapes of the graphs and indicate at which point choking first occurs
\lipsum[1-2]

% FIXME: wip
\subsection{Comparison to Theory} \label{sec:analysis/theory-comparison}

% TODO: compare limiting values of experimental \( \dot{m} \) and \( \frac{p_{throat}}{p_{atm}} \) with theoretical values
\lipsum[1-2]

% FIXME: wip
\section{Summary and Conclusions} \label{sec:summary}

% TODO: summary
\lipsum[1-2]

\section{References} \label{sec:references}

\printbibliography[heading=none]

\section{Appendices} \label{sec:appendices}

\subsection{Data} \label{sec:appendices/data}

\begin{table}[H]
    \centering
    \caption{Ambient Laboratory Conditions during Experiment}
    \begin{tabular}{|l|c|c|}
        \hline
        Quantity                  & Value                   & Uncertainty                     \\ \hline \hline
        Lab. Temperature, \( T \) & \qty{22}{\c}            & \(\pm \qty{0.5}{\c}\)           \\ \hline
        Lab. Pressure, \( P \)    & \qty{1025.1}{\hecto\Pa} & \( \pm \qty{0.05}{\hecto\Pa} \) \\ \hline
    \end{tabular}
    \label{tab:data/ambient-conditions}
\end{table}

\begin{table}[H]
    \centering
    \caption{Raw Measurements of Pressure Taps}
    \mintinline{bash}{./data/pressure_taps.csv} \\
    \begin{tabular}[c]{|c|c|c|}\hline
        \( P_{1} \left[\unit{\kPa}\right] \) & \( P_{2} \left[\unit{\kPa}\right] \) & \( P_{3} \left[\unit{\Pa}\right] \) \\ \hline \hline
        \csvreader[late after line = \\ \hline]{../data/pressure_taps.csv}{}{
            \csvcoli &
            \csvcolii &
            \csvcoliii
        }
    \end{tabular}
    \label{tab:data/pressure-taps}
\end{table}

\subsection{Code} \label{sec:appendices/code}

For the sake of brevity, only the code files that are key to the analysis are included below. However, in the spirit of completeness, the repository containing the complete data, source code, and notes for this report can be found at \href{https://www.github.com/vaisriv/enae464-lab03}{github:vaisriv/enae464-lab03}.

\begin{codeblock}
    \centering
    \caption{Index File}
    \mintinline{bash}{./src/index.py}
    \inputminted{python}{../src/index.py}
    \label{lst:code/index}
\end{codeblock}

% render all references
\nocite{*}

\end{document}
